% Source: https://artofproblemsolving.com/community/c3h1952183_geometry_problems__
% Stable identity: aops-topic-1952183-q02
The data are consistent with the power of the point $E$: $AE\cdot EB=20=CE\cdot ED$. Use coordinates with $E$ at the origin, $AB$ on the $x$-axis and $CD$ on the $y$-axis: $A=(-4,0)$, $B=(5,0)$, $C=(0,10)$, $D=(0,-2)$.

\begin{asy}
usepackage("amsmath");
size(5.5cm);
defaultpen(fontsize(10pt));
pair A=(-4,0),B=(5,0),C=(0,10),D=(0,-2),O=(0.5,4);real R=sqrt(36.25);
draw(circle(O,R));draw(A--B,linewidth(1));draw(C--D,linewidth(1));
draw(O--(0,4),red+linewidth(1));draw(O--(0.5,0),dashed);
draw((0,0)--(0.5,0)--(0.5,0.5)--(0,0.5)--cycle,gray);
draw((0,4)--(0.5,4)--(0.5,4.5)--(0,4.5)--cycle,gray);
dot(A);dot(B);dot(C);dot(D);dot(O);dot((0,0));
label("$A$",A,W);label("$B$",B,E);label("$C$",C,N);label("$D$",D,S);label("$O$",O,E);label("$E$",(0,0),SW);
label("$M$",(0,4),W);
label("\mbox{}",(min(currentpicture,user=true)+max(currentpicture,user=true))/2,black); // keeps the text after the figure black when LaTeX typesets the labels inline
\end{asy}

The centre $O$ lies on the perpendicular bisector of $AB$, the line $x=\dfrac{1}{2}$, and on the perpendicular bisector of $CD$, the line $y=4$. So $O=\left(\dfrac{1}{2},4\right)$ (and indeed $OA^2=4.5^2+4^2=36.25=0.5^2+6^2=OC^2$). The chord $CD$ lies on the line $x=0$, so the distance from $O$ to it is $\boxed{\dfrac{1}{2}}$.
